\documentclass[../../main.tex]{subfiles}
\begin{document}

\paragraph{【新課程】}
以下は平面内の問題である．
$O$，$A$，$B$，$C$は定点で，$A$，$B$，$C$は1直線上にないものとする．
\begin{enumerate}
  \item 点$P$が直線$AB$上にあるための必要十分条件は
    $\overrightarrow{OP}=a\overrightarrow{OA}+b\overrightarrow{OB}$，$a+b=1$
    ($a$，$b$は実数)と書けることである．これを証明せよ．
  \item 次の2条件をみたす実数$p$，$q$，$r$は$p=0$，$q=0$，$r=0$以外にないことを示せ．
    \begin{align}
      p\overrightarrow{OA}+q\overrightarrow{OB}+r\overrightarrow{OC}&=\vec{0} \quad (\text{零ベクトル}), \\
      p+q+r&=0
    \end{align}
  \item $Q$がこの平面上の点であって，
    $\overrightarrow{AQ}=x\overrightarrow{AB}+y\overrightarrow{AC}$ ($x$，$y$は実数)
    であるとき，
    $\overrightarrow{OQ}=l\overrightarrow{OA}+m\overrightarrow{OB}+n\overrightarrow{OC}$，
    $l+m+n=1$をみたす実数$l$，$m$，$n$は必ず存在し，
    しかもおのおのの値はただ1つに定まることを証明せよ．
\end{enumerate}

\paragraph{【旧課程】}
$\triangle ABC$の周上の2点$P$，$Q$を結ぶ線分$PQ$で3角形の面積を2等分する．
このような線分$PQ$の長さの最小値を求めよ．
ただし$BC$，$CA$，$AB$の長さをそれぞれ$a$，$b$，$c$とし，また$a>b>c$とする．

\end{document}
